\section{Conclusiones}

\subsection{Sostenibilidad y erradicación del trabajo manual}
\begin{indentedsubsection}
En estricta alineación con el Plan de Implementación (ProInnóvate), el presente proyecto marca el fin de la dependencia de procesos manuales, cuadernos físicos y cálculos repetitivos. La adopción del ecosistema Odoo no es un simple cambio de software, sino una inversión estratégica que permite a Perú Flores Travel Agency centralizar su base de datos. Al tener a los clientes y proveedores en una sola plataforma en la nube, la empresa asegura la integridad de su información, sentando bases reales y sólidas para escalar su volumen de pasajeros de manera ordenada.
\end{indentedsubsection}

\subsection{Agilidad comercial y reducción de tiempos reales}
\begin{indentedsubsection}
Tal como se definió en el Documento de Alcance, el uso del cotizador automatizado representa la mejora más tangible y rápida de este proyecto. Al automatizar el prorrateo de costos fijos (como movilidad y guías) y la simulación de ganancias variables (15\%, 20\%, 30\%), se proyecta cumplir la meta del Plan de Implementación de reducir los tiempos de cotización a menos de 15 minutos por prospecto. Esto otorga a la Gerencia General una ventaja competitiva inmediata: enviar propuestas profesionales en PDF al turista extranjero mucho más rápido que la competencia, sin tener que copiar, pegar ni sumar en Excel.
\end{indentedsubsection}

\subsection{Control financiero aterrizado y seguro}
\begin{indentedsubsection}
Una conclusión altamente práctica de esta implementación es la simplificación del control de pagos. El sistema se ha adaptado a la realidad exacta de la agencia, permitiendo registrar los adelantos (pagos ``a cuenta'') que ingresan por medios externos como Western Union, y calculando instantáneamente el saldo pendiente. Esta funcionalidad otorga a la gerencia total seguridad financiera y claridad visual antes de proceder con la compra de boletos o el pago a proveedores.
\end{indentedsubsection}

\subsection{Mejora continua y potenciamiento del ``Boca a Boca''}
\begin{indentedsubsection}
La configuración del módulo de Encuestas cumple un objetivo realista y de alto valor comercial: sistematizar el feedback del turista. Al automatizar el envío de formularios de satisfacción al finalizar el servicio, la agencia podrá recopilar testimonios que evalúen la calidad de los guías y el transporte. Esto no solo permite auditar el servicio, sino que genera valoraciones reales que impulsan la recomendación digital (``boca a boca''), factor fundamental para la captación de nuevos turistas en el mercado receptivo.
\end{indentedsubsection}
